\documentclass[11pt]{article}

% --------------------------------------------------
% Packages
% --------------------------------------------------
\usepackage[T1]{fontenc}
\usepackage{lmodern}
\usepackage{geometry}
\usepackage{microtype}
\usepackage{setspace}
\usepackage{amsmath,amssymb,amsthm,mathtools}
\usepackage{csquotes}
\usepackage{hyperref}
\usepackage{booktabs}

\geometry{margin=1in}
\setstretch{1.15}

% --------------------------------------------------
% Theorem environments
% --------------------------------------------------
\newtheorem{definition}{Definition}
\newtheorem{lemma}{Lemma}
\newtheorem{proposition}{Proposition}
\newtheorem{remark}{Remark}

% --------------------------------------------------
% Macros
% --------------------------------------------------
\newcommand{\Hist}{\mathcal{H}}
\newcommand{\Future}{\mathcal{F}}
\newcommand{\Adm}{\mathrm{Adm}}
\newcommand{\Reg}{\mathsf{Reg}}
\newcommand{\Val}{\mathsf{Val}}
\newcommand{\Do}{\mathrm{do}}

% --------------------------------------------------
% Title
% --------------------------------------------------
\title{Event-Nodes and Texture-Nodes:\\
Extending Causal Mechanistic Interpretability with Irreversible Structure}
\author{Flyxion}
\date{December 2025}

\begin{document}
\maketitle

% ============================================================
\begin{abstract}
Mechanistic interpretability has increasingly converged on a causal paradigm in
which neural networks are analyzed through interventions, counterfactuals, and
causal abstractions. While this framework successfully distinguishes correlation
from mechanism, it remains implicitly state-centered and treats irreversibility
as a contingent byproduct of information loss. This essay proposes a minimal,
formal extension of the causal–mechanistic framework by introducing two typed
classes of causal nodes: \emph{texture-nodes}, which carry instantaneous values,
and \emph{event-nodes}, which record regime selections that irreversibly constrain
future admissible behavior. Drawing an explicit analogy to procedural material
systems—where generator selection and parameter ratios play distinct roles—we
develop typing rules, a generalized causal tracing method, and graph rewrite
semantics that preserve causal rigor while accounting for history sensitivity,
structural irreversibility, and algorithmic regime switching. The result is a
unified account in which causal abstraction remains valid, but is explicitly
temporal, plural, and constraint-aware.
\end{abstract}

% ============================================================
\section{Introduction}
% ============================================================

Recent work in mechanistic interpretability has emphasized the necessity of
causal reasoning for understanding neural networks. Interventions, rather than
mere inspection, are taken as the criterion for distinguishing true mechanisms
from correlational artifacts. Within this paradigm, understanding a model means
reverse-engineering the algorithm it implements by demonstrating systematic
interventional correspondence between low-level components and high-level
computational structure.

Despite its success, this framework implicitly privileges instantaneous internal
states and treats history as relevant only insofar as it manifests in present
activations. Irreversibility appears primarily as epistemic degradation caused by
noise or ablation, rather than as a constitutive feature of mechanism itself. As
a result, several phenomena—commitment, collapse, regime switching, and
substitution failure—remain awkward to express within a purely state-based
causal graph.

This essay introduces a minimal extension that resolves these limitations without
abandoning causal rigor. The core move is a typed distinction between two kinds of
causal nodes: \emph{texture-nodes}, which carry values, and \emph{event-nodes},
which select generative regimes and irreversibly constrain future behavior. The
extension is minimal in the sense that it adds no new metaphysics, only new types
for causal variables and interventions.

% ============================================================
\section{From State-Based Causality to Typed Causal Nodes}
% ============================================================

Standard causal models represent systems as directed acyclic graphs whose nodes
carry variables updated by structural equations. Interventions replace the value
of a variable, and causal responsibility is assessed by downstream changes in
behavior.

We retain this structure but introduce a typing function that distinguishes
value-carrying nodes from constraint-carrying nodes.

\begin{definition}[Typed Causal Graph]
A typed causal model is a tuple
\[
\mathcal{M} = (G, \tau, \Val, \Reg, \Adm)
\]
where $G=(V,E)$ is a directed graph, $\tau:V\to\{\mathbf{T},\mathbf{E}\}$ assigns
each node a type, $\Val$ is the space of texture values, $\Reg$ is the space of
event regimes, and $\Adm$ maps realized histories to sets of admissible futures.
\end{definition}

Texture-nodes ($\mathbf{T}$) represent internal quantities whose causal role is
exhausted by their instantaneous value. Event-nodes ($\mathbf{E}$) represent
structural choices whose effects are not reducible to a value, but instead alter
which future states remain reachable.

This distinction parallels the difference between modifying parameters of a
procedural generator and switching the generator itself.

% ============================================================
\section{Texture-Nodes}
% ============================================================

A texture-node carries a value $X_v \in \Val_v$ and updates according to a
structural equation
\[
X_v := f_v(X_{\mathrm{Pa}(v)}, E_{\mathrm{Pa}(v)}).
\]

Interventions on texture-nodes are standard do-operations:
\[
\Do(v := x), \quad x \in \Val_v.
\]

Texture-nodes are fully replayable: their causal role can be tested by substituting
values and observing behavioral changes. Steering vectors, activation patching,
and difference-in-mean methods all operate at this level.

Importantly, texture values often arise from generator functions parameterized by
ratio-like quantities. These ratios behave as vector directions in parameter
space, but their meaning depends on the generator selected upstream.

% ============================================================
\section{Event-Nodes}
% ============================================================

An event-node carries a regime value $E_v \in \Reg_v$ that selects a generative
mode or structural configuration. Its update rule may depend on both texture and
event parents:
\[
E_v := g_v(X_{\mathrm{Pa}(v)}, E_{\mathrm{Pa}(v)}).
\]

Interventions on event-nodes take the form
\[
\Do(v := r), \quad r \in \Reg_v.
\]

Unlike texture interventions, event interventions need not be reversible. They
can permanently restrict the space of admissible futures.

\begin{definition}[Structural Irreversibility]
An intervention on an event-node $e$ is structurally irreversible at time $t$ if
\[
\Adm(\Hist_t) \neq \Adm(\Hist_t[R \leftarrow r']),
\]
where $\Hist_t$ is the realized history up to $t$.
\end{definition}

Irreversibility here is not informational but structural: even perfect replay
cannot restore the original counterfactual space.

% ============================================================
\section{Generator Selection and Parameter Ratios}
% ============================================================

Consider a texture-node generated by a procedural function $G_r$ selected by an
event-node with regime $r$:
\[
X = G_r(\rho(\theta), \ldots),
\]
where $\rho(\theta)$ denotes ratios of generator parameters.

Changing $\rho(\theta)$ modifies the realized texture without altering the
generator family. Changing $r$ selects a different generator altogether, producing
a qualitatively different space of possible outputs. Parameter ratios thus behave
as vectors only relative to a fixed regime.

This clarifies why some activation directions generalize while others fail under
context shifts: the underlying generator has changed.

% ============================================================
\section{Causal Tracing with Event-Nodes}
% ============================================================

Standard causal tracing corrupts an input and restores internal activations to
identify mediators. With typed nodes, two distinct restoration operations are
available.

\begin{itemize}
\item \textbf{Texture restoration:} $\Do(X_v := X_v^\star)$
\item \textbf{Event restoration:} $\Do(E_v := E_v^\star)$
\end{itemize}

Texture restoration measures recovery of output behavior. Event restoration
measures recovery of admissible futures. A node is a structural mediator if
restoring its regime reinstates counterfactual possibilities even when output
accuracy remains unstable.

This explains why some interventions preserve fluency while destroying coherence:
the texture is intact, but the regime has shifted.

% ============================================================
\section{Graph Rewrite Semantics}
% ============================================================

Event-nodes can be given explicit operational meaning as graph rewrite triggers.
The texture computation graph contains labeled rewrite sites. A rewrite rule has
the form
\[
(s, r) \Rightarrow G_s^{(r)},
\]
indicating that site $s$ is replaced by subgraph $G_s^{(r)}$ when regime $r$ is
active.

Rewrites must preserve interface types, ensuring that regime switching is
structural rather than destructive. Texture parameters live inside the selected
subgraph; event-nodes determine which subgraph exists.

This formalizes regime switching as algorithmic structure, not noise.

% ============================================================
\section{Causal Abstraction Revisited}
% ============================================================

Causal abstraction is traditionally validated by interventional equivalence.
Under the typed framework, abstraction requires two correspondences:

\begin{enumerate}
\item Texture correspondence: value-level interventions align.
\item Event correspondence: regime transitions and irreversibility align.
\end{enumerate}

An abstraction that preserves outputs but misrepresents regime structure is
behaviorally accurate but temporally incoherent. Multiple abstractions may remain
valid, but each carries an explicit scope and loss profile.

% ============================================================
\section{What Counts as Understanding}
% ============================================================

Understanding is no longer mere control or prediction. It is the ability to
reconstruct which regimes were selected, which futures were foreclosed, and how
value-level variation unfolds within those constraints.

A system may pass all texture-level tests while remaining opaque at the event
level. Such cases mark the boundary of state-based explanation, not its failure.

% ============================================================
\section{Conclusion}
% ============================================================

By distinguishing texture-nodes from event-nodes, the causal–mechanistic paradigm
is extended just enough to account for history sensitivity, irreversibility, and
regime selection without abandoning its core commitments. Interventions remain the
criterion of explanation, but their targets are now typed. Causal tracing becomes
structural as well as behavioral. Graph rewrite rules make explicit what was
previously implicit in algorithmic mixture.

The resulting framework preserves the strengths of causal interpretability while
making explicit the limits of state-only analysis. Mechanisms are not merely paths
through variables, but trajectories through constrained possibility spaces.

\end{document}
